% 核心结论：交流电桥的平衡
% 图原型：纯几何／对比。物理网络按本题节点、端口和符号构建。
\documentclass[tikz,border=7pt]{standalone}
\usepackage{fontspec,xeCJK,amsmath,amssymb}
\IfFontExistsTF{Times New Roman}{\setmainfont{Times New Roman}}{\setmainfont{TeX Gyre Termes}}
\IfFontExistsTF{Heiti SC}{\setCJKmainfont[AutoFakeBold=1.5]{Heiti SC}}{\setCJKmainfont[AutoFakeBold=1.5]{FandolHei-Regular}}
\usetikzlibrary{arrows.meta,calc,angles,quotes}
\definecolor{ink}{HTML}{243745}
\definecolor{blue}{HTML}{275DAB}
\definecolor{green}{HTML}{177C63}
\definecolor{red}{HTML}{BD4935}
\definecolor{purple}{HTML}{715393}
\begin{document}
\begin{tikzpicture}[font=\fontsize{9.5}{12}\selectfont,text=ink,
 v/.style={line width=1.1pt,-{Stealth[length=2.4mm,width=1.5mm]}},
 axis/.style={ink!65,line width=.45pt,-{Stealth[length=1.7mm,width=1.1mm]}},
 aux/.style={ink!40,densely dashed,line width=.5pt},
 lab/.style={fill=white,inner sep=2pt},
 card/.style={draw=blue!25,fill=blue!3,rounded corners=3pt,text width=6cm,align=left,inner sep=10pt}]
\path[use as bounding box] (-8,-3.35) rectangle (8,3.65);
\node[font=\fontsize{10.5}{13}\selectfont\bfseries,text width=15cm,align=center] at (0.00000,3.22000) {M6-020　交流电桥的平衡};
\begin{scope}[shift={(-4,0)}]
\draw[ink,thick] (-2.50000,1.50000) -- (-1.67000,1.50000);
\draw[ink,thick] (-0.83000,1.50000) -- (0.00000,1.50000);
\begin{scope}[shift={(-1.25000,1.50000)},rotate=0.0]
\draw[ink,thick,fill=white] (-0.42,-.14) rectangle (0.42,.14);
\end{scope}
\node[lab] at (-1.25000,1.92000) {$Z_1$};
\draw[ink,thick] (0.00000,1.50000) -- (0.83000,1.50000);
\draw[ink,thick] (1.67000,1.50000) -- (2.50000,1.50000);
\begin{scope}[shift={(1.25000,1.50000)},rotate=0.0]
\draw[ink,thick,fill=white] (-0.42,-.14) rectangle (0.42,.14);
\end{scope}
\node[lab] at (1.25000,1.92000) {$Z_2$};
\draw[ink,thick] (-2.50000,1.50000) -- (-1.51888,0.32265);
\draw[ink,thick] (-0.98112,-0.32265) -- (0.00000,-1.50000);
\begin{scope}[shift={(-1.25000,0.00000)},rotate=-50.19442890773481]
\draw[ink,thick,fill=white] (-0.42,-.14) rectangle (0.42,.14);
\end{scope}
\node[lab] at (-1.57265,-0.26888) {$Z_3$};
\draw[ink,thick] (0.00000,-1.50000) -- (0.98112,-0.32265);
\draw[ink,thick] (1.51888,0.32265) -- (2.50000,1.50000);
\begin{scope}[shift={(1.25000,0.00000)},rotate=50.19442890773481]
\draw[ink,thick,fill=white] (-0.42,-.14) rectangle (0.42,.14);
\end{scope}
\node[lab] at (1.57265,-0.26888) {$Z_4$};
\draw[ink,thick] (0.00000,1.50000) -- (0.00000,0.42000);
\draw[ink,thick] (0.00000,-0.42000) -- (0.00000,-1.50000);
\begin{scope}[shift={(0.00000,0.00000)},rotate=-90.0]
\draw[ink,thick,fill=white] (-0.42,-.14) rectangle (0.42,.14);
\end{scope}
\node[lab,anchor=east] at (-0.42000,0.00000) {$G$};
\fill[ink] (-2.50000,1.50000) circle (1.6pt);
\node[] at (-2.60000,1.80000) {$A$};
\fill[ink] (0.00000,1.50000) circle (1.6pt);
\node[] at (-0.10000,1.80000) {$B$};
\fill[ink] (2.50000,1.50000) circle (1.6pt);
\node[] at (2.40000,1.80000) {$C$};
\fill[ink] (0.00000,-1.50000) circle (1.6pt);
\node[] at (0.00000,-1.80000) {$D$};
\draw[ink,thick] (-2.50000,1.50000) -- (-2.90000,1.50000);
\draw[ink,thick] (-2.90000,1.50000) -- (-2.90000,-2.25000);
\draw[ink,thick] (-2.90000,-2.25000) -- (-0.40000,-2.25000);
\draw[ink,thick] (0,-2.25) circle(.4);
\node[] at (0.00000,-2.25000) {$\sim$};
\draw[ink,thick] (0.40000,-2.25000) -- (2.90000,-2.25000);
\draw[ink,thick] (2.90000,-2.25000) -- (2.90000,1.50000);
\draw[ink,thick] (2.90000,1.50000) -- (2.50000,1.50000);
\end{scope}
\begin{scope}[shift={(4,0)}]
\node[card] at (0.00000,1.15000) {$Z_1Z_4=Z_2Z_3$};
\end{scope}
\begin{scope}[shift={(4,0)}]
\node[card] at (0.00000,-1.15000) {$R_x=100\;\Omega,\quad C_x=10\;\mu F$};
\end{scope}
\node[text width=15cm,align=center] at (0.00000,-3.00000) {桥臂AB/BC与AD/DC构成两条分压路径，B、D等势给出交叉乘积条件。};

\end{tikzpicture}\end{document}